% ---------------------------------------------------------------------------
% The short version's introduction. Poses the proposal's question, says plainly
% what the proposal got right and what it got wrong, and promises only what the
% five sections after it deliver. Numbers are quoted from the complete paper and
% carry the same registry keys.
% ---------------------------------------------------------------------------

\section{Introduction}
\label{s:sec:intro}

An observability bill is one of the few software costs that can grow faster than the
system it measures, and the published cases are startling: one company's monitoring
invoice reached \$65M a year, run up, by the account its own engineers give, because
nobody was watching infrastructure spend~\citep{orosz2023}. The obvious response is to
move to something cheaper, and for metrics the cheap option is well established ---
Prometheus is free, widely deployed, and its query language is a de facto standard. Teams
nonetheless stay, and when they explain why, the reason is rarely the data. It is the
dashboards.

A mature observability estate is thousands of saved queries: panels, alerting rules,
service-level objectives, runbook links. Moving vendors means rewriting all of them in
another language, by hand, without a way to tell whether a rewrite preserved what the
original meant. That is a switching cost, and the question worth asking is where it comes
from. If it is incidental --- different keywords, different punctuation --- then a
translator removes it and the lock-in was never real. If it follows from how the two
languages are \emph{designed}, then no amount of engineering removes it, and the query
language is a retention mechanism whether or not anyone intended it as one.

The question is tractable because query languages are formal languages. Rich treats them
as such (§48.1), and the book's classification --- regular, context-free, decidable ---
is the frame this paper works in~\citep{rich2008}. If the two languages are context-free,
they have grammars; if their grammars are unambiguous, each query has one parse tree; if
each construct of one has a meaning-preserving counterpart in the other, a translator
exists; and where a construct does not, the obstruction can be exhibited as a theorem
instead of reported as a limitation.

\paragraph{What was proposed, and what the measurement did to it.}
This work was proposed as a single question --- is the lock-in a syntactic accident or a
semantic necessity? --- with a plan attached: formalise both languages as
$(\text{grammar},\text{semantics})$ pairs, define what it means for a translation between
them to be sound, build a syntax-directed transpiler through a shared intermediate
representation, and measure on real dashboards how much of each language actually
translates. That plan was carried out as written, and three things happened to it.

The first is a correction the proposal could not have anticipated, and it changed the
definition of the central term. Soundness cannot be an equation between two queries
evaluated on one database, because the two systems do not receive the same database: a
Prometheus counter is a running total and a Datadog counter a sequence of per-interval
deltas. Soundness here is therefore stated relative to an \emph{ingestion mapping}
$\phi$, and every theorem names the class of databases it holds on
(\S\ref{s:sec:framework}).

The second is that the expected answer did not arrive. The proposal took for granted that
some constructs would turn out to be semantically inexpressible, and the paper reserved a
column for them. In the forward direction that column emptied of everything we can
demonstrate, bar $14$ occurrences left undecided: of the $8{,}770$ Datadog metric query
occurrences in our corpus, $96.6\%$ translate, and none of the shortfall is shown to be
beyond PromQL's reach. Most of the
corrections that emptied it were to our own earlier claims rather than to the vendor's
behaviour. Forward lock-in, on this evidence, is work somebody has not done.

The third is that the necessity was real, and in the direction the proposal treated as
secondary. Going back --- PromQL to Datadog --- a tag filter must become a glob, globs
are a proper subclass of the regular languages, and the refusals there are theorems
rather than unfinished work. That asymmetry is this paper's result, and it is the reason
this version is organised around the reverse direction rather than the forward one.
LogQL and Splunk SPL were named in the proposal as stretch targets and remain stretch
targets; no grammar or translator for either appears here.

\paragraph{Scope.}
We study \emph{forward portability}: once a team dual-writes or switches, a translated
query must mean the same thing as the original on the data each system receives from then
on. Migrating stored history is a different problem and is out of scope. We measure the
query layer only. The main translation direction is Datadog metric queries into PromQL,
because that is the direction teams actually take; the reverse direction is evidence
about tag filtering drawn from the other half of the corpus, with no translator built for
it, and the theorems of \S\ref{s:sec:regular} are about what any translator in that
direction could achieve rather than about one we wrote.

\paragraph{What this version carries.}
This is the short version of a longer paper, and it carries one of that paper's results
in full. \S\ref{s:sec:background} places the question among the four literatures it
belongs to, the last of which predicts the shape of the answer.
\S\ref{s:sec:framework} fixes the framework: what the systems store, what a
query denotes, when two denotations agree, and what a translator may do when it cannot
proceed. \S\ref{s:sec:forward} reports the forward direction compactly --- the grammars
exist and admit deterministic one-token-lookahead parsers, the transpiler is
syntax-directed through an intermediate representation, and the coverage it achieves is
the number above. \S\ref{s:sec:regular} is the core: globs, what they provably cannot
say, a decision procedure for whether a regular expression has a glob at all, and what
the corpus turns out to contain. \S\ref{s:sec:asymmetry} draws the asymmetry and names
the finding this version leaves out. \S\ref{s:sec:conclusion} answers the proposal's
question. Proofs of the forward-direction results, the grammar conflict analysis, the
ingestion mapping in full, the evaluation apparatus and the experiments put to live
accounts are in the complete paper and are cited here by name rather than reproduced.
